% Replaces the figure float labelled fig:pca-repr (accepted main.tex, Appendix: the float that inputs
% figures/pca-walking.tikz). Label unchanged, so every \Cref keeps working.
% Needs: figures/pca-walking.tikz  <- projects/pmt_camera_ready/writing/rainbow_figures_20261102/tikz/pca-walking.tikz
%        data/{wave_hybrid,patch_pts_hybrid,state_pts_hybrid,state_boundary_hybrid}_walking-walking_downstairs-walking.txt
%        <- the same record's data/
% Splice: HAR validation 18 / 0 / 18 (Walking / Walking Downstairs / the same Walking recording), tokenizer
% tokens and last-block memory states, position-centred (record README, "Position check").
% Replaced-float marker: the manuscript colours \revtab captions green and changed wording
% (\rev) blue until \revfinaltrue.
\providecommand{\revtab}[1]{#1}
\providecommand{\rev}[1]{#1}
\begin{figure}[!ht]
\centering
\resizebox{\linewidth}{!}{%
% fig:pca-repr: PCA, validation splice 18 / 0 / 18 (Walking, Walking Downstairs, the same Walking sequence).
% One colour ramp per segment (Dark2-A / Dark2-B / Dark2-C, light = start of the segment, dark = its end),
% marker = segment; dashed = linear boundary fitted to the plotted states. Columns seg / src / segpos are
% appended to the accepted format (export_rainbow.py). Needs pgfplots >= 1.13 (per-plot colormap name).
\begin{tikzpicture}
\begin{groupplot}[
  group style={group name={myplot}, group size=3 by 1, horizontal sep=1.25cm},
  footnotesize,
  width=5cm,
  height=5cm,
  table/col sep=comma,
  ylabel near ticks,
  ylabel shift=-5pt,
  ticklabel style={font=\tiny},
  colormap={seg1}{rgb255(0cm)=(118,197,173); rgb255(4cm)=(27,158,119); rgb255(6.5cm)=(20,119,89)},
  colormap={seg2}{rgb255(0cm)=(232,159,103); rgb255(4cm)=(217,95,2); rgb255(7.5cm)=(141,62,1)},
  colormap={seg3}{rgb255(0cm)=(152,148,198); rgb255(2.5cm)=(117,112,179); rgb255(5.5cm)=(82,78,125)},
  point meta min=0,
  point meta max=1,
  shaded/.style={scatter, only marks, point meta=\thisrow{segpos}, forget plot,
    scatter/use mapped color={draw=black!65, fill=mapped color}, mark options={line width=0.25pt}},
  wave/.style={mesh, thick, point meta=\thisrow{segpos}, forget plot},
]

\nextgroupplot[xlabel=Time, ylabel=Amplitude, xmin=0, xmax=1, ymin=-5, ymax=5]
  \addplot[wave, colormap name=seg1, restrict expr to domain={\thisrow{seg}}{0:0}] table[x=norm_pos, y=ch5] {data/wave_hybrid_walking-walking_downstairs-walking.txt};
  \addplot[wave, colormap name=seg2, restrict expr to domain={\thisrow{seg}}{1:1}] table[x=norm_pos, y=ch5] {data/wave_hybrid_walking-walking_downstairs-walking.txt};
  \addplot[wave, colormap name=seg3, restrict expr to domain={\thisrow{seg}}{2:2}] table[x=norm_pos, y=ch5] {data/wave_hybrid_walking-walking_downstairs-walking.txt};

\nextgroupplot[xlabel=PC 1, ylabel=PC 2]
  \addplot[shaded, mark=*, mark size=1.2pt, colormap name=seg1, restrict expr to domain={\thisrow{seg}}{0:0}] table[x=dim0, y=dim1] {data/patch_pts_hybrid_walking-walking_downstairs-walking.txt};
  \addplot[shaded, mark=square*, mark size=1.2pt, colormap name=seg2, restrict expr to domain={\thisrow{seg}}{1:1}] table[x=dim0, y=dim1] {data/patch_pts_hybrid_walking-walking_downstairs-walking.txt};
  \addplot[shaded, mark=triangle*, mark size=1.2pt, colormap name=seg3, restrict expr to domain={\thisrow{seg}}{2:2}] table[x=dim0, y=dim1] {data/patch_pts_hybrid_walking-walking_downstairs-walking.txt};

\nextgroupplot[
  xlabel=PC 1,
  ylabel=PC 2,
  legend to name=pcalegend,
  legend style={legend columns=4, font=\tiny, draw=none, fill=none, /tikz/every even column/.append style={column sep=6pt}},
]
  \addplot[shaded, mark=*, mark size=1.8pt, colormap name=seg1, restrict expr to domain={\thisrow{seg}}{0:0}] table[x=dim0, y=dim1] {data/state_pts_hybrid_walking-walking_downstairs-walking.txt};
  \addplot[shaded, mark=square*, mark size=1.8pt, colormap name=seg2, restrict expr to domain={\thisrow{seg}}{1:1}] table[x=dim0, y=dim1] {data/state_pts_hybrid_walking-walking_downstairs-walking.txt};
  \addplot[shaded, mark=triangle*, mark size=1.8pt, colormap name=seg3, restrict expr to domain={\thisrow{seg}}{2:2}] table[x=dim0, y=dim1] {data/state_pts_hybrid_walking-walking_downstairs-walking.txt};
  \addlegendimage{only marks, mark=*, mark size=1.8pt, draw=black!65, fill=Dark2-A, mark options={line width=0.25pt}}
  \addlegendentry{Walking (seg.~1)}
  \addlegendimage{only marks, mark=square*, mark size=1.8pt, draw=black!65, fill=Dark2-B, mark options={line width=0.25pt}}
  \addlegendentry{Downstairs (seg.~2)}
  \addlegendimage{only marks, mark=triangle*, mark size=1.8pt, draw=black!65, fill=Dark2-C, mark options={line width=0.25pt}}
  \addlegendentry{Walking (seg.~3)}
  \addplot[black!60, dashed, thick] table[x=x, y=y] {data/state_boundary_hybrid_walking-walking_downstairs-walking.txt};
  \addlegendentry{Fitted linear boundary}
\end{groupplot}
% One legend row, centred over the token and state panels (it applies to both).
\node[anchor=south] at ($(myplot c2r1.north west)!0.5!(myplot c3r1.north east) + (0,3pt)$) {\pgfplotslegendfromname{pcalegend}};
\end{tikzpicture}%
%
}
\caption{\revtab{(Left) Double spliced HAR waveform. (Middle) PCA of representations (\rev{tokenizer} tokens) of the signal. (Right) PCA of memory-states. \rev{The splice joins two validation recordings: Walking, Walking Downstairs and the same Walking recording again. Tokens and memory states are shown after removing the mean representation at each position over the validation set. Colour marks the segment and shade the position within it (light = start); the dashed line is a linear boundary fitted to the memory states.}}}
\label{fig:pca-repr}
\end{figure}
